\documentclass[11pt,a4paper]{article}

\usepackage{iftex}
\ifPDFTeX
  \usepackage[utf8]{inputenc}
  \usepackage[T1]{fontenc}
  \usepackage{lmodern}
\else
  \usepackage{fontspec}
\fi
\usepackage[portuguese]{babel}
\usepackage{microtype}
\usepackage[margin=2.5cm]{geometry}
\usepackage{booktabs}
\usepackage{tabularx}
\usepackage{array}
\usepackage{fancyvrb}
\usepackage[font=small,labelfont=bf]{caption}
\usepackage{url}
\usepackage{placeins}
\usepackage{pgfplots}
\pgfplotsset{compat=1.18}
\usetikzlibrary{babel}
\usepackage[hidelinks]{hyperref}

\renewcommand{\arraystretch}{1.15}
\newcolumntype{L}{>{\raggedright\arraybackslash}X}
\newcolumntype{P}[1]{>{\raggedright\arraybackslash}p{#1}}

\title{Relatório de desempenho\\[0.4em]\large Resposta médica de emergência comunitária descentralizada\\ISIA 2026/27, trabalho prático 1, tópico O}
\author{Rodrigo Dias (202407281) \and Francisco Machado (202403514)}
\date{FCUP, novembro de 2026}

\begin{document}

\maketitle

\section{Introdução}

Este relatório avalia o sistema multiagente em SPADE que desenvolvemos para o tópico O, resposta médica de emergência comunitária sem despachante central. Segue o que o enunciado pede para o relatório de desempenho: resultados em pelo menos três cenários de teste contrastantes, análise das métricas, trade-offs e limitações. Os agentes, os formatos de mensagem e os protocolos estão descritos em \path{docs/} e nos slides; aqui resumimos só o necessário para ler os resultados.

A cidade está dividida em células, e o armário de DAE de cada célula é um hub: o agente que verifica os alertas da célula, coordena os casos e mantém a reputação das fontes. Para declarar um incidente, o hub funde a evidência dos sensores dos cidadãos e dos relatos das testemunhas, e pesa cada fonte pela confiança que tem nela, calculada com o modelo de reputação Beta (Jøsang e Ismail, 2002). Voluntários, DAE e ambulâncias são alocados por Contract Net (Smith, 1980; FIPA, 2002), com o ETA de cada um como proposta. A testemunha recebe no telemóvel uma guia de RCP passo a passo, ativada logo na triagem, antes de o caso ser declarado. Cada hub replica os seus casos num sucessor fixo, que os assume quando deixa de receber os heartbeats do hub.

\section{Método}

\subsection{Cenários}

Definimos cinco cenários a partir de uma cidade base (C1), mais do que os três pedidos, porque cada mecanismo do sistema só é posto à prova numa situação própria. Cada cenário altera poucas coisas face a C1 e responde a uma pergunta (Tabela~\ref{tab:cenarios}). Os parâmetros estão em \path{scenarios/}.

\begin{table}[htbp]
\centering
\small
\caption{Cenários de teste.}\label{tab:cenarios}
\begin{tabularx}{\textwidth}{P{2.2cm}LL}
\toprule
Cenário & O que muda face a C1 & Pergunta \\
\midrule
C1 base urbana & Cidade de 3 $\times$ 3 km com 10 hubs, 30 voluntários e 40 dispositivos; sensores fiáveis, rede sem perdas, 3 incidentes por hora & O fluxo completo funciona? Com que frequência chega um voluntário antes da ambulância? \\
C2 ruído & Falsos positivos dos sensores $\times$5, quatro sensores em deriva, testemunhas que leem mal a cena & A verificação mantém a precisão? Que política de alerta deve o sensor seguir? \\
C3 pico & Oito incidentes em 10 minutos num raio de 400 m & A alocação evita conflitos entre casos? Quanta redundância compensa? \\
C4 periferia & Metade dos voluntários, metade dos DAE fechados, ambulâncias fora da área, 5~\% de perda, duas zonas sem rede, falha do hub 2 aos 1200 s & A guia e a alocação compensam a distância? Quanto custa recuperar de um hub que falha? \\
C5 adversarial & A cidade de C1 com duas pessoas que relatam colapsos falsos, quatro por hora cada & A reputação trava relatos falsos repetidos, e a que custo? \\
\bottomrule
\end{tabularx}
\end{table}

\subsection{Configurações comparadas}

Em cada cenário comparamos o sistema tal como foi desenhado (``nosso'') com a baseline B1 e com variantes que desligam ou alteram um mecanismo de cada vez.

A B1 é uma difusão ingénua: declara o caso no primeiro alerta, chama todos os voluntários até ao segundo anel de células, adjudica aos primeiros que respondem e não usa reputação (\path{docs/03} §4.2). Como não faz triagem, também não ativa a guia da testemunha. O B0 é um oráculo centralizado, calculado offline sobre o mesmo traço de incidentes: sabe onde e quando ocorre cada colapso e onde estão todos os voluntários. Não é atingível e serve só de referência para a métrica M9.

As variantes são: sem guia da testemunha (C1, C3 e C4); sem reputação (C2 e C5); as políticas de alerta do sensor AlertFirst, que alerta logo, e ConfirmFirst, que pergunta primeiro ao portador, em vez da adaptativa, que só alerta logo perante sinais de P1 (C2); redundância de $k = 1$ e $k = 3$ voluntários de RCP por P1, em vez de $k = 2$ (C3); e perda de 10~\% e 20~\% das mensagens nas ligações dos telemóveis (C1). Com a B1 em todos os cenários, são 21 configurações. A variante sem guia também serve para isolar a coordenação: como nem ela nem a B1 têm guia, a diferença entre as duas mede o efeito da verificação, do Contract Net e da reputação.

\subsection{Métricas}

As métricas foram definidas na análise Gaia (\path{docs/01}, §8.1) e são calculadas a partir do event log de cada run (\path{src/cemr/metrics/compute.py}). As de tempo contam desde o colapso, usam a mediana dentro de cada run e só incluem os incidentes detetáveis, isto é, com um dispositivo ou uma testemunha. A Tabela~\ref{tab:metricas} lista as que usamos neste relatório. O M4 (tempo até ao DAE) e o M13 (equidade entre centro e periferia) também foram calculados, mas têm poucos valores por configuração (o M4 tem $n = 1$ em C2) e não os discutimos; estão em \path{results/summary.md}.

\begin{table}[htbp]
\centering
\small
\caption{Métricas usadas.}\label{tab:metricas}
\begin{tabularx}{\textwidth}{P{4.2cm}L}
\toprule
Métrica & Definição \\
\midrule
M1 latência de verificação & Tempo até à declaração do caso (P1 e P2) \\
M2 verificação & Precisão dos casos declarados; recall sobre os P1/P2 detetáveis ou reportados; falsos despachos por run (caso declarado onde não há doente) \\
M3 tempo até à RCP & Tempo até ao início da RCP (P1) e percentagem iniciada pela testemunha \\
M5 tempo até à ambulância & Tempo até a ambulância chegar (P1 e P2) \\
M6 proxy de sobrevivência & Fórmula de Larsen et al.\ (1993) aplicada aos P1 desfibrilháveis; serve só para comparar configurações \\
M7 cobertura antes da EMS & Percentagem de P1 com voluntário (ou DAE) no local antes da ambulância \\
M8 mensagens & Mensagens por candidato a caso, sem as de fundo (registos, gossip de reputação e heartbeats) \\
M9 distância ao oráculo & Atraso do primeiro voluntário face ao B0 \\
M10 carga dos voluntários & Alertas por voluntário, índice de Jain das tarefas e minutos gastos em falsos alarmes \\
M11 exposição & Agentes que recebem a localização exata do doente, por caso \\
M12 recuperação & Tempo até o sucessor assumir os casos de um hub que falhou; casos órfãos \\
\bottomrule
\end{tabularx}
\end{table}

\subsection{Procedimento e análise estatística}\label{sec:estat}

Cada configuração correu com as seeds 1 a 10, num total de 210 runs de 60 minutos simulados, com o relógio acelerado 10 vezes. Os 210 runs são válidos (o relógio simulado avançou como devia) e foram todos feitos sobre o mesmo commit (\texttt{de13290}). A seed fixa a cidade, as posições iniciais e o traço de incidentes, por isso as configurações de um mesmo cenário veem exatamente os mesmos incidentes.

Para cada configuração apresentamos a média sobre as seeds com o intervalo de confiança (IC) a 95~\% da t de Student. As comparações são emparelhadas por seed: para cada seed calculamos a diferença variante $-$ nosso, e reportamos a média dessas diferenças, o seu IC a 95~\% e em quantas seeds o nosso sistema foi melhor. Só tiramos conclusões de diferenças cujo IC exclui 0.

Uma métrica fica indefinida num run sem eventos para a medir, por exemplo o M3 num run em que ninguém chegou a fazer RCP. O $n$ de cada média é o número de seeds em que a métrica existe. Isto torna otimistas algumas médias da B1: em C1, a B1 só tem M3 em 7 das 10 seeds, porque nas outras três ninguém fez RCP.

O verificador de regras (\path{src/cemr/metrics/check.py}) relê o event log de cada run e conta as violações das regras organizacionais da análise Gaia (um só dono por caso, uma reserva por DAE, divulgação da localização só a quem precisa, entre outras) e da regra DUP (nenhum doente com dois casos declarados). Os resultados reproduzem-se com:

\begin{Verbatim}[frame=single,fontsize=\footnotesize]
python experiments/run_matrix.py --matrix main --seeds 1-10 --workers 4
python experiments/analyse.py        # summary, paired, misses e figuras em results/
python -m cemr.metrics.check runs/main/<run> -v
\end{Verbatim}

\section{Resultados}\label{sec:resultados}

Os números desta secção vêm de \path{results/summary.md} (médias), \path{results/paired.md} (diferenças emparelhadas) e \path{results/misses.md} (incidentes falhados por motivo). A Tabela~\ref{tab:b1} compara o nosso sistema com a B1 nos cinco cenários, e a Figura~\ref{fig:m3} mostra o tempo até à RCP.

\begin{table}[htbp]
\centering
\footnotesize
\setlength{\tabcolsep}{4pt}
\caption{Nosso sistema e B1 nos cinco cenários: médias sobre as seeds, com IC a 95~\%. As médias de M3 da B1 usam 7, 3, 10, 9 e 6 seeds (secção~\ref{sec:estat}).}\label{tab:b1}
\begin{tabular}{llccccc}
\toprule
Métrica & & C1 & C2 & C3 & C4 & C5 \\
\midrule
M1 declaração (s) & Nosso & 33,5 $\pm$ 10,7 & 40,6 $\pm$ 13,8 & 49,5 $\pm$ 3,5 & 41,3 $\pm$ 9,0 & 39,9 $\pm$ 12,9 \\
 & B1 & 26,4 $\pm$ 8,6 & 30,5 $\pm$ 10,7 & 39,1 $\pm$ 4,7 & 27,7 $\pm$ 10,0 & 30,6 $\pm$ 10,6 \\
\addlinespace
M2 precisão (\%) & Nosso & 100 & 23 $\pm$ 14 & 100 & 100 & 28 $\pm$ 15 \\
 & B1 & 100 & 25 $\pm$ 11 & 100 & 100 & 15 $\pm$ 8 \\
\addlinespace
M2 recall, detetáveis (\%) & Nosso & 93 $\pm$ 12 & 76 $\pm$ 33 & 80 $\pm$ 11 & 96 $\pm$ 7 & 83 $\pm$ 26 \\
 & B1 & 93 $\pm$ 12 & 83 $\pm$ 26 & 82 $\pm$ 10 & 96 $\pm$ 7 & 83 $\pm$ 26 \\
\addlinespace
M2 falsos despachos/run & Nosso & 0 & 4,7 $\pm$ 2,0 & 0 & 0 & 3,0 $\pm$ 0,7 \\
 & B1 & 0 & 9,4 $\pm$ 2,1 & 0 & 0 & 7,1 $\pm$ 1,1 \\
\addlinespace
M3 tempo até à RCP (s) & Nosso & 52 $\pm$ 12 & 66 $\pm$ 26 & 75 $\pm$ 25 & 164 $\pm$ 108 & 83 $\pm$ 64 \\
 & B1 & 307 $\pm$ 41 & 271 $\pm$ 134 & 362 $\pm$ 81 & 383 $\pm$ 83 & 280 $\pm$ 61 \\
\addlinespace
M6 sobrevivência (\%) & Nosso & 41,8 $\pm$ 4,7 & 36,2 $\pm$ 14,9 & 24,8 $\pm$ 10,8 & 16,1 $\pm$ 6,8 & 36,7 $\pm$ 15,2 \\
 & B1 & 31,3 $\pm$ 4,0 & 26,9 $\pm$ 11,5 & 13,9 $\pm$ 7,0 & 9,1 $\pm$ 5,4 & 27,9 $\pm$ 12,0 \\
\addlinespace
M7 voluntário antes da EMS (\%) & Nosso & 82 $\pm$ 17 & 56 $\pm$ 41 & 75 $\pm$ 16 & 91 $\pm$ 11 & 56 $\pm$ 41 \\
 & B1 & 56 $\pm$ 32 & 26 $\pm$ 33 & 49 $\pm$ 15 & 70 $\pm$ 23 & 48 $\pm$ 39 \\
\addlinespace
M8 mensagens/candidato & Nosso & 166 $\pm$ 10 & 105 $\pm$ 7 & 218 $\pm$ 16 & 131 $\pm$ 15 & 111 $\pm$ 12 \\
 & B1 & 127 $\pm$ 5 & 93 $\pm$ 8 & 146 $\pm$ 11 & 94 $\pm$ 10 & 110 $\pm$ 6 \\
\addlinespace
M10 alertas/voluntário & Nosso & 1,3 $\pm$ 0,6 & 3,0 $\pm$ 0,9 & 4,5 $\pm$ 0,5 & 1,9 $\pm$ 0,6 & 2,3 $\pm$ 0,7 \\
 & B1 & 2,4 $\pm$ 1,0 & 5,9 $\pm$ 0,1 & 5,8 $\pm$ 0,3 & 3,1 $\pm$ 1,0 & 5,4 $\pm$ 0,5 \\
\bottomrule
\end{tabular}
\end{table}

\begin{figure}[htbp]
\centering
\begin{tikzpicture}
\begin{axis}[
  ybar,
  bar width=10pt,
  width=0.85\textwidth,
  height=6cm,
  symbolic x coords={C1,C2,C3,C4,C5},
  xtick=data,
  ymin=0,
  ylabel={Tempo até à RCP (s)},
  enlarge x limits=0.15,
  legend style={at={(0.5,1.03)}, anchor=south, legend columns=3, draw=none,
    /tikz/every even column/.append style={column sep=0.6cm}},
  legend image code/.code={\draw[#1] (0cm,-0.1cm) rectangle (0.3cm,0.15cm);},
  error bars/y dir=both,
  error bars/y explicit,
  error bars/error bar style={black},
]
\addplot[fill=black!75, draw=black] coordinates {
  (C1,52.2) +- (0,12.4) (C2,65.8) +- (0,26.4) (C3,75.1) +- (0,25.2)
  (C4,164.4) +- (0,107.5) (C5,82.8) +- (0,64.0)};
\addplot[fill=black!35, draw=black] coordinates {
  (C1,229.8) +- (0,51.8) (C3,339.3) +- (0,36.6) (C4,411.1) +- (0,52.2)};
\addplot[fill=white, draw=black] coordinates {
  (C1,306.5) +- (0,41.3) (C2,270.8) +- (0,134.3) (C3,361.6) +- (0,81.0)
  (C4,383.1) +- (0,83.0) (C5,279.5) +- (0,61.0)};
\legend{Nosso, Sem guia, B1}
\end{axis}
\end{tikzpicture}
\caption{M3, tempo até à RCP nos P1: média sobre as seeds com IC a 95~\%. A variante sem guia só correu em C1, C3 e C4.}\label{fig:m3}
\end{figure}

Em todos os cenários, o nosso sistema põe alguém a fazer RCP 2 a 6 vezes mais cedo do que a B1, e o proxy de sobrevivência fica 7 a 11 pontos acima. As diferenças emparelhadas de M3 excluem 0 em C1, C3, C4 e C5 (em C1, +258 $\pm$ 39 s, com o nosso melhor nas 7 seeds em que as duas configurações têm RCP). Em C2 só duas seeds têm RCP nas duas configurações, e o IC é demasiado largo para uma conclusão. As diferenças de M6 excluem 0 nos cinco cenários.

A Figura~\ref{fig:m3} mostra também de onde vem este ganho. Sem guia, o tempo até à RCP do nosso sistema aproxima-se do da B1: em C3 e C4 a diferença entre os dois não se distingue de zero, e em C1 fica 55 s abaixo. A maior parte do ganho na RCP vem, portanto, da guia da testemunha, e não da coordenação entre agentes. A coordenação nota-se noutras métricas: menos alertas por voluntário (M10) em todos os cenários e, nos cenários com falsos alarmes, metade ou menos dos falsos despachos. Também põe um voluntário no local antes da ambulância em mais casos (M7), mas essa diferença só exclui 0 em C3 e, na comparação B1 $-$ sem guia, em C4; em C1, C2 e C5 o IC inclui 0.

Em sentido contrário, o nosso sistema declara os casos mais tarde em todos os cenários (M1), e em nenhuma seed de nenhum cenário declarou mais depressa do que a B1. A secção~\ref{sec:tradeoffs} discute este custo.

A Tabela~\ref{tab:efeitos} reúne as diferenças emparelhadas das variantes, que as subsecções seguintes discutem cenário a cenário.

\begin{table}[htbp]
\centering
\footnotesize
\caption{Efeito de cada variante: diferença emparelhada variante $-$ nosso, média $\pm$ IC a 95~\%, e em quantas seeds o nosso sistema foi melhor. Nas linhas ``B1 $-$ sem guia'' a referência é a variante sem guia, e a última coluna conta as seeds em que ela foi melhor do que a B1. ``p.p.'' são pontos percentuais.}\label{tab:efeitos}
\begin{tabularx}{\textwidth}{llLrc}
\toprule
Cenário & Variante & Métrica & $\Delta$ & Melhor \\
\midrule
C1 & Sem guia & M3 tempo até à RCP (s) & +178 $\pm$ 51 & 10/10 \\
   & Sem guia & M6 sobrevivência (p.p.) & $-$7,0 $\pm$ 2,6 & 10/10 \\
   & B1 $-$ sem guia & M3 tempo até à RCP (s) & +55 $\pm$ 51 & 6/7 \\
   & B1 $-$ sem guia & M10 alertas por voluntário & +1,1 $\pm$ 0,4 & 10/10 \\
   & Perda 10~\% & M3 tempo até à RCP (s) & +123 $\pm$ 102 & 10/10 \\
   & Perda 10~\% & M3 RCP pela testemunha (p.p.) & $-$40 $\pm$ 33 & 5/10 \\
   & Perda 20~\% & M6 sobrevivência (p.p.) & $-$21,8 $\pm$ 12,3 & 10/10 \\
   & Perda 20~\% & M2 recall, reportados (p.p.) & $-$16 $\pm$ 23 & 3/10 \\
\addlinespace
C2 & Sem reputação & M2 falsos despachos/run & +1,1 $\pm$ 0,8 & 6/10 \\
   & Sem reputação & M10 minutos em falsos alarmes & +11,8 $\pm$ 9,0 & 8/10 \\
   & AlertFirst & M2 falsos despachos/run & $-$0,1 $\pm$ 0,2 & 0/10 \\
   & AlertFirst & M8 mensagens/candidato & $-$12 $\pm$ 4 & 0/10 \\
   & ConfirmFirst & M1 declaração (s) & +3,9 $\pm$ 10,5 & 5/7 \\
   & ConfirmFirst & M8 mensagens/candidato & +16 $\pm$ 8 & 10/10 \\
\addlinespace
C3 & Sem guia & M3 tempo até à RCP (s) & +264 $\pm$ 34 & 10/10 \\
   & B1 $-$ sem guia & M3 tempo até à RCP (s) & +22 $\pm$ 68 & 6/10 \\
   & B1 $-$ sem guia & M7 voluntário antes da EMS (p.p.) & $-$25 $\pm$ 17 & 7/10 \\
   & $k = 1$ & M11 agentes com a localização & $-$1,1 $\pm$ 0,4 & 0/10 \\
   & $k = 1$ & M10 índice de Jain & $-$0,2 $\pm$ 0,1 & 9/10 \\
   & $k = 3$ & M11 agentes com a localização & +1,7 $\pm$ 0,8 & 9/10 \\
   & $k = 3$ & M9 distância ao oráculo (s) & $-$19 $\pm$ 49 & 4/10 \\
\addlinespace
C4 & Sem guia & M3 tempo até à RCP (s) & +247 $\pm$ 97 & 9/10 \\
   & B1 $-$ sem guia & M3 tempo até à RCP (s) & $-$18 $\pm$ 117 & 3/9 \\
   & B1 $-$ sem guia & M7 voluntário antes da EMS (p.p.) & $-$25 $\pm$ 23 & 5/10 \\
   & B1 $-$ sem guia & M7 DAE antes da EMS (p.p.) & $-$34 $\pm$ 31 & 6/10 \\
\addlinespace
C5 & Sem reputação & M2 falsos despachos/run & +3,9 $\pm$ 1,4 & 10/10 \\
   & Sem reputação & M2 precisão (p.p.) & $-$12,5 $\pm$ 8,4 & 8/10 \\
   & Sem reputação & M10 minutos em falsos alarmes & +24,1 $\pm$ 14,7 & 9/10 \\
   & Sem reputação & M3 tempo até à RCP (s) & $-$0,3 $\pm$ 36,8 & 3/7 \\
\bottomrule
\end{tabularx}
\end{table}

\FloatBarrier
\subsection{C1 base urbana}

Na cidade base, o fluxo completo funciona sem erros de verificação: precisão e recall sobre os incidentes reportados de 100~\%, nenhum falso despacho e nenhuma violação de regras. Com a guia, é a testemunha que começa 95~\% das RCP; sem ela, a RCP começa 178 s mais tarde. A coordenação sozinha (B1 $-$ sem guia) encurta a RCP em 55 $\pm$ 51 s, um IC que exclui 0 por pouco, e reduz os alertas por voluntário.

A perda de mensagens nas ligações dos telemóveis é onde o sistema mais se degrada. Com 10~\% de perda, a RCP começa 123 s mais tarde e a percentagem iniciada pela testemunha cai de 95~\% para 55~\%. Com 20~\%, o proxy de sobrevivência cai para 20,0~\%, abaixo dos 31,3~\% da B1 sem perda (esta comparação não é emparelhada), e o recall sobre os reportados cai para 84~\%, com um IC que inclui 0. Uma explicação provável está no código: a ativação da guia (P05) é enviada uma só vez a cada testemunha (\path{activate_guidance} em \path{src/cemr/agents/case_verify.py}) e não é reenviada quando se perde. Não a confirmámos nos event logs. As duas variantes com perda somam 6 das 16 violações de regras da matriz: cinco casos duplicados entre células (DUP) e uma testemunha que disse que o doente não respondia e não chegou a receber a guia (regra OL3).

\subsection{C2 ruído e falsos alarmes}

Com sensores ruidosos, a precisão fica entre 22~\% e 25~\% em todas as configurações e não difere entre o nosso sistema e a B1 ($\Delta$ = +1,8 $\pm$ 12,1). Parte da explicação é estrutural: há cerca de dois P1/P2 por run (23 nas dez seeds), e com tão poucos casos verdadeiros alguns falsos despachos chegam para baixar a precisão. A verificação nota-se no número de falsos despachos, que cai para metade (4,7 contra 9,4 por run; $\Delta$ = +4,7 $\pm$ 1,4, melhor nas 10 seeds). Quase toda a diferença está nos alertas só de sensores (2,9 contra 7,4 por run), que são os que a corroboração filtra. Os falsos despachos com testemunha quase não mudam (1,8 contra 2,0), porque uma testemunha que lê mal a cena passa a verificação como relato humano fiável.

O custo está no recall. A B1 tem mais 7,4 $\pm$ 17,1 pontos sobre os detetáveis, um IC que inclui 0. No total das seeds, o nosso sistema perdeu dois P1/P2 que a B1 apanhou: um descartado como P3 na triagem e outro fundido com o caso de outro doente.

As políticas AlertFirst e ConfirmFirst não diferem da adaptativa em precisão, recall, falsos despachos ou M1; só nas mensagens, com a AlertFirst a gastar menos 12 por candidato e a ConfirmFirst mais 16. Com o hub a corroborar os alertas, perguntar primeiro ao portador não trouxe nada mensurável. A reputação tem aqui um efeito pequeno (mais 1,1 falsos despachos por run sem ela, melhor em 6 de 10 seeds): cada sensor em deriva dá cerca de três falsos positivos por hora, e em 60 minutos são poucos casos para a confiança nele descer.

\subsection{C3 pico com contenção}

No pico, o recall sobre os detetáveis é de 80~\%, igual ao da B1. Dos 94 P1/P2 das dez seeds, 9 eram indetetáveis, 14 não chegaram a ser reportados e 3 foram fundidos com o caso de outro doente. Estes 3 aparecem também na B1, por isso vêm da fusão de alertas próximos e não da verificação. As regras de exclusividade (um compromisso por voluntário, uma reserva por DAE) não foram violadas em nenhum dos 50 runs de C3.

A redundância $k$ não muda o tempo até à RCP (75 s com $k = 1$, 2 ou 3), porque com a guia é a testemunha que começa 86 a 88~\% das RCP. Muda os custos. Com $k = 3$, a localização exata do doente chega a mais 1,7 agentes por caso; com $k = 1$, chega a menos 1,1, mas as tarefas concentram-se em menos voluntários e o índice de Jain cai 0,2. O ganho de $k = 3$ na distância ao oráculo (19 s) tem um IC que inclui 0.

Sem guia, a coordenação não encurta a RCP face à B1 ($\Delta$ = +22 $\pm$ 68 s), mas põe um voluntário no local antes da ambulância em mais 25 pontos percentuais dos P1.

\subsection{C4 periferia degradada}

Na periferia, a ambulância demora em média 760 s, cerca de 13 minutos. O nosso sistema começa a RCP em 164 s, contra 383 s com a B1, e o proxy de sobrevivência sobe de 9,1~\% para 16,1~\%. A guia explica a diferença na RCP: sem ela, o M3 sobe para 411 s, e a coordenação sozinha não o encurta ($\Delta$ = $-$18 $\pm$ 117 s). O que a coordenação traz é chegar antes da ambulância: mais 25 pontos percentuais de P1 com voluntário e mais 34 com DAE, ambos com IC que exclui 0 por pouco.

O hub 2 falha aos 1200 s em todos os runs de C4. O sucessor assume os casos em 21,1 $\pm$ 1,7 s, sem casos órfãos nos 30 runs. Este tempo é quase todo fixado pelos parâmetros (lease de três heartbeats de 5 s e um heartbeat de margem, 20 s), por isso o M12 mostra que o failover funciona, mas não mede uma propriedade emergente.

O nosso sistema não teve violações de regras em C4. A B1 teve 3 e a variante sem guia 6, todas casos duplicados entre células. O aviso entre hubs vizinhos que resolve estes duplicados (P26) não é esperado antes da declaração, e com 5~\% de perda e zonas sem rede chega tarde ou perde-se mais vezes. Não analisámos porque é que a variante sem guia duplica mais casos do que o sistema completo.

\subsection{C5 fontes adversárias}

As duas fontes adversárias descrevem uma pessoa inconsciente que respira (P2), um relato que, sem reputação, passa a verificação como relato humano fiável. Com reputação, os falsos despachos caem de 6,9 para 3,0 por run (melhor nas 10 seeds), a precisão sobe de 15,8~\% para 28,3~\% e os voluntários gastam metade do tempo em falsos alarmes (22,9 contra 47,0 minutos por run). Não encontrámos custo nos casos verdadeiros: o tempo até à RCP e o recall são os mesmos com e sem reputação.

Os 3,0 falsos despachos que restam vêm da forma como o modelo Beta aprende. A confiança numa fonte só desce quando um caso em que ela participou é fechado como falso alarme (P21, \path{docs/03} §4.5), por isso os primeiros relatos de cada fonte adversária ainda passam.

\section{Trade-offs}\label{sec:tradeoffs}

Verificar custa tempo. O nosso sistema declara os casos 7 a 14 s depois da B1 (diferença emparelhada de M1 entre 7,0 $\pm$ 2,7 s em C1 e 13,6 $\pm$ 7,7 s em C4), e esse atraso passa para a ambulância, que com a B1 chega 7 a 12 s mais cedo em C1, C2 e C3. Não passa para a RCP, porque a guia é ativada na triagem, antes da declaração. O que se ganha em troca depende do cenário. Com falsos alarmes (C2 e C5), os falsos despachos caem para metade ou menos. Sem falsos alarmes (C1, C3 e C4), a precisão é de 100~\% nas duas configurações, e a verificação só custa tempo.

O nosso sistema troca mais mensagens por candidato do que a B1 (mais 30~\% em C1, 49~\% em C3 e 39~\% em C4), mas cada voluntário recebe menos alertas, de menos 22~\% em C3 a menos 57~\% em C5. Preferimos esta troca, porque o recurso escasso é a atenção dos voluntários e não a rede. As mensagens dos casos são, aliás, uma parte pequena do total: as de fundo (registos, gossip de reputação e heartbeats) são entre 14~600 e 16~200 por run, 97~\% do total em C1, e não dependem do número de incidentes. São o custo fixo da descentralização: a descoberta de voluntários e testemunhas, a partilha de reputação entre hubs e a deteção de hubs que falham.

Na redundância, mais voluntários por P1 espalham as tarefas e aproximam o primeiro voluntário do oráculo, mas enviam a localização exata do doente a mais pessoas. Com a guia ativa, $k$ não muda o tempo até à RCP, e o ganho de $k = 3$ em M9 não se distingue de zero. Mantivemos $k = 2$ como compromisso entre a exposição e a concentração das tarefas, sabendo que os dados não mostram um ganho de tempo com mais redundância.

\section{Evolução do desenho}\label{sec:evolucao}

Três alterações ao desenho vieram da leitura de matrizes anteriores. Cada uma foi medida nas mesmas seeds antes e depois (Tabela~\ref{tab:evolucao}; \path{results/evolution.md}, \path{docs/03} §7.1). Entre duas medições pode haver mais do que uma alteração, por isso cada linha mede o efeito de tudo o que mudou entre elas.

\begin{table}[htbp]
\centering
\small
\caption{Alterações ao desenho medidas nas mesmas seeds, antes e depois.}\label{tab:evolucao}
\begin{tabularx}{\textwidth}{LP{4.2cm}P{4.6cm}}
\toprule
Alteração & Medida & Antes $\rightarrow$ depois \\
\midrule
Fundir alertas por área de 100 m, em vez de por célula & C3, recall detetáveis, seeds 1-5 & 64,0 $\rightarrow$ 77,5~\% (+13,5 $\pm$ 13,3) \\
Repetir a triagem que contradiz o relato & C3, recall detetáveis, seeds 1-10 & 75,8 $\rightarrow$ 80,3~\% (+4,5 $\pm$ 4,2) \\
A mesma fonte conta uma só vez por candidato & C5, falsos despachos por run, seeds 1-10 & 3,0 $\rightarrow$ 3,0 (sem efeito) \\
\bottomrule
\end{tabularx}
\end{table}

A fusão por célula juntava num só caso dois doentes da mesma célula. A triagem repetida acabou com os P1/P2 descartados como P3 em C3 (de 3 para 0, segundo o diagnóstico de \path{src/cemr/metrics/misses.py}). A terceira alteração não teve efeito em C5, porque aí os falsos despachos vêm de testemunhas e cada uma envia um só relato por incidente. Mantivemo-la porque é a fusão de evidência correta (fontes independentes), e não por um ganho medido.

\section{Limitações}\label{sec:limitacoes}

O ambiente é sintético. A cidade é uma grelha de 3 $\times$ 3 km com deslocações de Manhattan, e o comportamento das pessoas (aceitar uma chamada, responder à triagem, desistir da RCP) segue modelos com parâmetros escolhidos por nós, sem calibração empírica. A ambulância tem 240 s de ativação e desloca-se a 30 km/h; não encontrámos uma fonte para os tempos de resposta urbanos em Portugal, e os valores absolutos de M5 e M6 dependem desta escolha.

Não simulamos o caminho do 112. Os incidentes sem dispositivo nem testemunha contam como falhados, e as métricas de tempo só cobrem os detetáveis. A comparação com a B1 é feita sobre os mesmos incidentes, mas não com o sistema de emergência real. O M6 usa o modelo de Larsen et al.\ (1993), feito para paragens testemunhadas em fibrilhação ventricular; serve para comparar configurações e não é uma previsão clínica.

A amostra é pequena. Com 10 seeds de 60 minutos há poucos P1 por run (16 no total das seeds em C2 e C5), e várias métricas têm IC largos ou poucos valores. Runs mais longos dariam também mais tempo à reputação para aprender com sensores em deriva.

Ficam dois defeitos conhecidos. Os casos duplicados entre células (15 das 16 violações de regras da matriz, nenhuma na configuração principal) aparecem quando o aviso P26 entre hubs vizinhos se perde ou chega depois da declaração. E a ativação da guia não é reenviada quando se perde, o que provavelmente explica a quebra da RCP pela testemunha com perda de mensagens.

Por fim, dentro de um processo o SPADE entrega as mensagens em memória, sem passar pelo servidor XMPP (\path{spikes/README.md}). A latência, a perda e as zonas sem rede vêm do nosso modelo de rede e não de uma rede real.

\section{Conclusões}\label{sec:conclusoes}

O sistema cumpre o objetivo principal do tópico: chega ao doente mais cedo do que a difusão ingénua nos cinco cenários, sem despachante central. A principal lição da avaliação é que a maior parte do ganho no tempo até à RCP vem da guia da testemunha, um mecanismo simples, e não da coordenação entre agentes. A coordenação contribui com menos alertas por voluntário e, com fontes pouco fiáveis, metade ou menos dos falsos despachos, à custa de 7 a 14 s na declaração. Nos cenários de pico e de periferia (C3 e C4) leva ainda mais vezes um voluntário ao local antes da ambulância, e em C4 também o DAE. Os resultados negativos são a precisão baixa com sensores muito ruidosos, igual à da B1, e a degradação com perda de mensagens.

Como trabalho futuro, propomos esperar pela resposta P26 antes de declarar um caso perto da fronteira de uma célula, reenviar a ativação da guia quando se perde, simular o caminho do 112 e correr runs mais longos para a reputação ter mais casos de onde aprender.

\section*{Referências}

\begin{list}{}{\setlength{\leftmargin}{1.5em}\setlength{\itemindent}{-1.5em}}
\item FIPA (2002). \emph{FIPA Contract Net Interaction Protocol Specification} (SC00029H).
\item Jøsang, A., \& Ismail, R. (2002). The Beta Reputation System. \emph{15th Bled Electronic Commerce Conference}.
\item Larsen, M. P., Eisenberg, M. S., Cummins, R. O., \& Hallstrom, A. P. (1993). Predicting survival from out-of-hospital cardiac arrest: a graphic model. \emph{Annals of Emergency Medicine}, 22(11), 1652--1658.
\item Smith, R. G. (1980). The Contract Net Protocol. \emph{IEEE Transactions on Computers}, C-29(12), 1104--1113.
\end{list}

\end{document}
